\section{Supplementary: Method Decision Log and Stage 3 Pipeline Contract}
\label{sec:suppl:nextsteps}

\subsection{Multi-point substrate-scale scoping curve}
\label{sec:suppl:nextsteps:substrate-curve}
The substrate-extension probe at $\SubstrateProbeCpCropSize^3$ (suppl.\ \S\ref{sec:suppl:exp:substrate-probe-cprime}) establishes one point of substrate-extent contrast against the $\SubstrateProbeCropSize^3$ probe at fixed simulation pitch ($n_{\text{grid}}=768$). Two points define a contrast, not a curve. A scoping curve sampling additional substrate edges (e.g., $96^3$, $128^3$, and beyond) at the same fixed pitch, under the same five-seed protocol, is the natural next step before any verb stronger than ``contrast at $(\SubstrateProbeCropSize^3, \SubstrateProbeCpCropSize^3)$'' can be supported. The compute budget required scales as the cube of the substrate edge for the truth-$\rho$ field at fixed pitch and is the binding constraint on extending to $\geq 96^3$; a budget expansion is the operational prerequisite. The pre-committed seed-variance budget should also be re-fitted to the empirical $\SubstrateProbeCpResnetSigma$ value (a factor of $\SubstrateProbeCpSeedFactor$ above the design budget; suppl.\ \S\ref{sec:suppl:exp:substrate-probe-cprime}) before a curve is drawn, so that the per-point CI width is consistent with the observed seed sensitivity~\cite{damour2020underspecification}.

\subsection{Stage 3: Feedback identification from reconstructed $\rho$}
\label{sec:suppl:nextsteps:stage3}
Stage 3 trains a four-way classifier on 3D crops of the reconstructed $\rho/\bar\rho$ field over $\{\text{P1}, \text{P2}, \text{P3}, \text{P4}\}$, with target accuracy $> 85\%$. The data-pipeline contract is fixed: input is a $(C{=}1, D, H, W)$ cube of reconstructed overdensity at native simulation resolution, label is the originating \texttt{physics\_id}. The architectural choice (3D CNN vs.\ 3D ViT, crop dimensions, augmentation policy) is deferred to a Stage 3 design pass and is gated on a $P_F$-passing reconstruction; the structural-$P_F$ gap of [D-39] is upstream of the feedback-discrimination claim.

\subsection{Method decision log}
\label{sec:suppl:nextsteps:decisions}
The methodology of \S\ref{sec:method} and the evaluation contract of \S\ref{sec:experiments} are the consolidated form of a longer decision sequence. Tab.~\ref{tab:decision-log} summarizes the scientific entries; full traceability lives in the project's experiment LEDGER.

\begin{table}[t]
    \centering
    \renewcommand{\arraystretch}{1.15}
    \setlength{\tabcolsep}{4pt}
    \begin{tabular}{l|p{0.74\columnwidth}}
    ID & Decision \\
    \hline
    {[D-01]} & Tepper-Garc\'ia 4th-order Voigt; small-$|x|$ Taylor branch added 2026-05-04 \\
    {[D-02]} & Bounded-activation physics layers for $\rho$, $T$, $X_{HI}$, $v_{\text{pec}}$ \\
    {[D-06]} & RSD-convolved integrator (Stage 2a$\to$2b lift); first-half = redshift-space target \\
    {[D-11]/[D-34]} & Mean-flux soft anchor $\langle F\rangle_{\text{obs}} = \AnchorMeanF \pm \AnchorMeanFsigma$ at $z=0.3$ (Kirkman+ 2007); cost-survey runs trained against retracted value $0.877$, kept under anchor-invariance \\
    {[D-13]} & Stage 2b gates: $P_F(k_\parallel)$, Pearson $\xi_{\hat\rho,\rho}(2\,h^{-1}\,\text{Mpc})$, flux-PDF KS on $F\in[\PDFwinLo,\PDFwinHi]$ \\
    {[D-14]} & Production schedule (per-tier max\_steps, microbatch, accum) \\
    {[D-19]} & Smoke-pass criteria: gradient finite, descent $\leq 0.85$, mean\_F within $5\%$ \\
    {[D-21]} & Two-pass surrogate gradient for the masked mean-flux constraint \\
    {[D-23]} & Tier-aware microbatch / step schedule with linear VRAM model \\
    {[D-24]} & log1p+cap+mask Ly$\alpha$ forest loss; $\tau_{\max}=10$ LOCKED 2026-05-04 \\
    {[D-35]} & Cost-survey post-hoc rescale read as anchor-preview, retained under [D-39] only as upper-bound optimism \\
    {[D-38]} & log-space supervision is the methods contribution; cap and mask retained as defense-in-depth \\
    {[D-39]} & \texttt{pub-t1} corrected-anchor full-schedule bundle falsifies the rescale-preview; $P_F$ is the structurally binding gate; Tier 4 conditional on Wrinkle-1 \\
    {[D-40]} & Saturation-aware $P_F$ band loss tested on $P_1$ (12.5k-step counterfactual): eval $|\Delta P_F/P_F|=\PFdfortyEval$ ($+37\%$ vs.\ baseline), KS $=\KSdfortyEval$ (gate broken); integrated-residual formulation admits constant-prediction degeneracy. \#4.1 ranking re-pivoted to per-pixel FGPA-tail prior. \\
    \end{tabular}
    \caption{Methodological decision log. One-line summary per scientific entry; LEDGER \S3 carries full derivations, sources, and verification status. [D-39] supersedes the verdict-frame implicit in earlier decision entries: the corrected-anchor full-schedule run is recorded as a calibration-vs-structure falsification, not a thresholded PASS.}
    \label{tab:decision-log}
\end{table}
